\chapter{Introduction}
\label{chap:introduction}

An evaluation is a lens through which we view the capabilities of a model.
A core challenge is then to construct evaluations that identify metrics that are both informative and unbiased, advancing our scientific understanding
of the current state of the art and make informed decisions about how to move forward.
In this thesis, we share three lessons regarding the evaluation of language models,
particularly for coding, mathematics, and logical reasoning.

\section{Lesson 1: Focusing on the Right Metrics Without Bias is Crucial}

The first lesson we will examine is that focusing on the right metrics without bias is crucial. In science,
almost anything can be measured, but not everything is meaningful. The first question we face when designing any
evaluation is what we should measure. 

In late 2023, language models like GPT-4 were just starting to become proficient at writing Python code.
When humans learn programming, they not only learn to write code, but also gain the ability to read code, understand it,
simulate its execution behavior, debug it, among a suite of other programming-related skills. At the time, there was little knowledge about 
whether language models learned to do programming in the same ways that humans did. Especially because the coding data that language models were
trained on primarily consisted of GitHub code, it was unclear whether these models could generalize to capabilities beyond code generation.
Motivated by this, we designed CRUXEval, a benchmark asking language models to simulate (forwards and reverse) the execution of very simple 
Python programs consisting of 3-13 lines with no complex control flow or calculations.

When we released CRUXEval, many people in the community were skeptical about the design and utility of such a benchmark. 
After all, a simple way to ace the forward execution subset of CRUXEval is to simply run the code in a Python interpreter.
However, we still saw value in this benchmark for several reasons. First, as is evident today, language models that can effectively 
program should have a good understanding and mental model of code. This may have been overlooked when simple single-turn benchmarks
like HumanEval were the standards for coding evaluations, because real-world programming requires a much wider set of skills beyond
single-shot code generation. Second, the benchmark allows us to scientifically understand new phenomena of coding language models.
For example, by showing that language models did have some ability to perform both forward and reverse execution, we saw early signs that
language models did have some degree of generalization despite just training on raw code tokens.

Beyond choosing the right metric, we must also account for potential biases in our evaluation design before arriving at conclusions.
As an example, in IneqMath, we introduced a benchmark for evaluating the ability of language models to solve inequality problems. 
When we measured the accuracy of Grok 3 mini, we saw that it was able to get the answer correct $72\%$ of the time, indicating a high level of proficiency at solving inequalities.
However, upon closer examination, when changing the metric from evaluating at the final answer to evaluating the entire reasoning trace and solution, we found that the
accuracy dropped to $6\%$, a stark difference that leads to an entire different conclusion. Arriving at the first conclusion misses the fact that models
may be able to correctly guess the final answer without providing a correct proof.

Another example showing the importance of accounting for bias and confounding factors is our LiveCodeBench evaluation. 
Prior to LiveCodeBench, coding evaluations primarily relied on static benchmarks like HumanEval, fixed datasets that were released once and 
never updated. However, this fails to account for data contamination: language models are often trained on data from the internet, which 
could potentially include these benchmarks. While many models are trained with best-effort decontamination processes, it is not possible to 
completely remove all traces of these benchmarks from the training set. Using time-controlled evaluations, we eliminate this bias. 
As evidence of this, DeepSeek-Instruct-33B has a published data cutoff date of September 2023. When evaluated on problems between May-August 2023, it had an average 
accuracy of around 50\%. However, when evaluated on problems between September 2023 and February 2024, the average accuracy dropped to around 20\%, significantly lower.
Similarly, the accuracy of DeepSeek-Base-33B (same cutoff date) plummeted from a pass@1 of 60\% on May 2023 LeetCode problems to around 0\% on September 2023 problems.

Later in the thesis, we will go into much more detail about these examples and others, showing that focusing on the
right metrics can lead to new insights and fresh perspectives on the capabilities of language models.
Once we have an evaluation and metric built, the next step is to use it to analyze the capabilities of language models.
Through in-depth analysis, we can often discover interesting phenomena that are not obvious from just
looking at the final scores, leading to the next lesson.

\section{Lesson 2: Looking Beyond the Score Reveals Hidden Nuances}

Often, the most interesting scientific insights can only be discovered after looking beyond surface-level metrics 
with deeper analysis. In CRUXEval, by changing the measurement metric, we were able to reveal that language models 
had some ability to perform execution. However, deeper analysis revealed even more insights about the relationship between 
different coding capabilities. The most striking finding was when we looked at the relationship between HumanEval scores and
CRUXEval scores across a suite of models. While the forward execution accuracy was highly correlated with the reverse execution
accuracy, the HumanEval score showed a different trend. For base models, performance on HumanEval and CRUXEval are correlated.
However, distilled models (WizardCoder, Phind, Phi) broke this trend, significantly beating their base models at HumanEval but
not CRUXEval. This reveals a key limit of fine-tuning: while distillation on NL-to-code samples can greatly boost performance on
code generation, it does not necessarily improve a model's underlying understanding of code.

Another illustration of this is when we developed LINC, a neurosymbolic method for logical reasoning.
The task is that given set of natural language premises and a conclusion, one must determine if the conclusion
is true, false, or unknown based on the premises. The standard method of doing this was via chain-of-thought (CoT)
prompting, having a language model reason through the premises to make a deduction about the validity of the conclusion.
Our method, LINC, first formalized everything into a first-order logic program and then executed a first-order logic prover
to determine the answer. On most of the evaluations, LINC outperformed CoT, which might lead one to believe it is
a superior method for performing logical reasoning. However, in one case, on the FOLIO benchmark 
using GPT-4 as the base model, this was not the case: LINC only scored $73\%$ vs. CoT's $75\%$. 

After this result, we conducted a series of thorough analyses to understand why LINC had underperformed in this
specific case. First, we did a qualitative analysis of the errors of both approaches, and discovered they were very 
different. Errors in LINC occurred when the models failed to properly capture implicit and explicit information that was 
in the natural language statements (failure modes impossible in CoT), while errors in CoT occurred from making incorrect natural language logical deductions
and failing to find the right reasoning chain (failure modes impossible in LINC). Upon further quantitative analysis, we found
that LINC had worse recall but better precision on True/False predictions, and that LINC and CoT mispredicted on different examples.
These nuances were not apparent from just looking at the performance scores and only came to light after
this deeper analysis.

Designing the right metric and carrying out the right analysis is only the first step.
Once an evaluation reveals these nuances, it naturally raises further questions about what
capabilities are missing, which failure modes matter most, and what kinds of systems should be
built next. In this way, evaluations not only describe the current state of models, but also
influence the trajectory of future research. This is the final lesson: what we measure
shapes what we build.

\section{Lesson 3: What We Measure Shapes What We Build}

Evaluations should always be designed to be forward-looking. Once a benchmark is released, the nuances
it uncovers often become a focal point for the relevant scientific community, influencing both what problems
researchers measure and what model-building labs optimize for. \benchmark, for example, influenced the community in
several ways.

First, many later benchmark efforts explicitly built upon the idea of measuring capabilities beyond code 
generation. For example, CRUXEval-X \citep{xu2024cruxevalx} extends our benchmark to measure execution in 19 different
programming languages. REval \citep{chen2024reval} evaluates runtime state prediction in addition to input/output
prediction. CodeMind \citep{liu2024codemind} introduces additional code reasoning tasks such as specification reasoning and
dynamic semantics reasoning. The analysis we did in \benchmark{} also inspired several
diagnostic and probing studies to understand the scientific properties of code language models. For example, in our
counterfeit conundrum work \citep{gu2024counterfeit}, we studied how language models reason about subtly incorrect code, using execution as a
probe to reveal that language models often simulated the execution of incorrect programs as if they were correct.
In SeqCoBench \citep{maveli2025seqcobench}, the authors tested whether code LLMs recognized functionally equivalent code
by applying semantics-preserving transformations and found that LLMs have a shallow understanding of semantics.
Rethinking Reflection in Pre-Training \citep{shah2025rethinkingreflection} uses \benchmark{} as a code reasoning task to test
whether pre-trained checkpoints can recover from incorrect prior reasoning about program inputs and outputs.

Second, many follow-up works have leveraged the idea of code execution prediction to better improve code 
language models. For example, SemCoder \citep{ding2024semcoder} introduces a new strategy called monologue reasoning in order to teach code language models
to reason about comprehensive semantics. CodeI/O \citep{li2025codei} found that structured training on
input-output predictions not only improves performance on \benchmark{}, but also more broadly on other reasoning tasks.
Later, Self-Execution Simulation \citep{maimon2026selfexecution} showed that combining supervised fine-tuning on execution traces
with reinforcement learning allows models to iteratively self-fix by simulating test execution.
Several months after its release, \benchmark{} also became a standard target for measuring the reasoning abilities of new
coding models. For example, DeepSeek-Coder V2 \citep{deepseekai2024deepseekcoderv2}, Qwen 2.5 \citep{qwen2024qwen25coder},
Qwen 3 \citep{qwen2025qwen3technical}, Kimi K2 \citep{moonshotai2025kimik2}, StarCoder 2 \citep{starcoder2},
and Code World Model \citep{faircodegen2025cwm} all cite and report performance on \benchmark{}.

\section{Works Covered in This Thesis}

The technical chapters of this thesis synthesize and discuss a collection of research projects on evaluating
language models for coding, mathematics, and reasoning. Chapter~\ref{chap:right-metrics} introduces \benchmark{}
\citep{gu2024cruxeval}, SlopCodeBench \citep{orlanski2026slopcodebench}, IneqMath \citep{lu2025ineqmath},
the Counterfeit Conundrum \citep{gu2024counterfeit}, and LiveCodeBench \citep{jain2024livecodebench}.
Chapter~\ref{chap:beyond-score} returns to several of these works and also introduces LINC \citep{olausson2023linc}.
Chapter~\ref{chap:measure-shapes} additionally introduces Vibe Code Bench \citep{vibe-code-bench}, LeanDojo \citep{yang2023leandojo}, and
ProofOptimizer \citep{gu2026proofoptimizer}. Chapter~\ref{chap:position-paper} closes with the position paper
\textit{Future Research and Challenges Remain Towards AI for Software Engineering} \citep{gu2025positionaiswe}.

Needless to say, these projects were not done alone. Many of the works discussed in this thesis were collaborations,
and the full author lists in the cited papers reflect the people who helped shape the ideas, build the
systems, run the experiments, and write the original papers. Throughout the thesis, ``we'' refers to the research teams behind the corresponding
works, and I am grateful to all of these collaborators for their contributions to all the content in this thesis.
